\documentclass[11pt]{article}
\usepackage[T1]{fontenc}
\usepackage[utf8]{inputenc}
\usepackage{lmodern}
\usepackage[margin=1in]{geometry}
\usepackage{amsmath,amssymb,amsthm,mathtools}
\usepackage{enumitem}
\usepackage[hidelinks]{hyperref}

\newtheorem{theorem}{Theorem}[section]
\newtheorem{proposition}[theorem]{Proposition}
\newtheorem{definition}[theorem]{Definition}
\newtheorem{lemma}[theorem]{Lemma}
\newtheorem{corollary}[theorem]{Corollary}
\newtheorem{remark}[theorem]{Remark}

\title{Generic Odd-Pfaffian Reconstruction from Commutator-Word Moments}
\author{Luca Blanchi}
\date{June 2026}

\begin{document}
\maketitle

\begin{abstract}
Let $k=\mathbb F_p$, $p\neq2$, and let

\[
\beta:\Lambda^2_k V\to W
\]

be an alternating tensor with $\dim V=2r+1$ odd. Equivalently, $\beta$ is an $m=\dim W$-dimensional linear system of alternating forms on $V$. For $[\lambda]\in\mathbb P(W^*)$, write

\[
A_\lambda=\lambda\circ\beta.
\]

The first nontrivial degeneracy locus is the corank-$\ge3$ Pfaffian locus

\[
Z_\beta=\{[\lambda]\in\mathbb P(W^*):\dim\ker A_\lambda\ge3\},
\]

cut out by the submaximal Pfaffians of the universal odd skew-symmetric matrix.

We prove that, on a natural nonempty Zariski open locus, pointwise Fourier-tagged scalar-extension commutator-word moments determine $\beta$ up to pseudo-isometry. The argument has four parts. First, the Fourier-tagged moments recover the rank of $A_\lambda$ for every $\lambda$ over every finite extension, hence recover the reduced Pfaffian degeneracy scheme $Z_\beta$. Under the explicit hypothesis that the submaximal Pfaffian ideal $I_\beta$ is saturated and geometrically radical, this recovers the ideal $I_\beta$. Second, by the Buchsbaum--Eisenbud structure theorem for codimension-three Gorenstein ideals, $I_\beta$ recovers the universal skew-symmetric matrix $M_\beta$ up to left-right equivalence of the middle map in the minimal free resolution. Third, if the common self-adjoint algebra

\[
\operatorname{Adj}^{+}(\beta)
=
\{T\in\operatorname{End}(V):T^{\mathsf T}A_\lambda=A_\lambda T\text{ for all }\lambda\}
\]

is only the scalar algebra $kI$, then left-right equivalence between alternating matrices upgrades to congruence, hence to pseudo-isometry of the tensor. Fourth, the Pfaffian hypotheses and the scalar self-adjoint condition hold on a nonempty Zariski open subset for

\[
\dim V=2r+1\ge5,\qquad 4\le \dim W\le10.
\]

Thus, for this range, Fourier-tagged MVMT gives a generic reconstruction theorem beyond alternating pencils. For the associated special class-two exponent-$p$ groups, it gives a generic word-moment reconstruction theorem for the commutator tensor, and hence for the group, on the surjective radical-free part of the open locus.

The theorem is not a universal classification of all alternating tensors and is not an optimized isomorphism algorithm. Its contribution is a generic reconstruction mechanism: commutator-word moments recover a Pfaffian codimension-three degeneracy scheme, and on a rigid open locus that scheme determines the tensor.

\end{abstract}

\section*{Declaration of generative AI and AI-assisted technologies in the writing process}
This article belongs to the Presentation theory section. Generative AI, including ChatGPT by \mbox{OpenAI}, was used to assist with mathematical drafting, formalization, review, and editing. The note may be preliminary, may have been only partially reviewed, and in some cases may be only AI-reviewed.

\section{Introduction}

Let $G$ be a special finite $p$-group of nilpotency class two and exponent $p$, with $p\neq2$. Then

\[
V=G/Z(G),\qquad W=G'=Z(G)
\]

are $\mathbb F_p$-vector spaces, and the commutator defines an alternating tensor

\[
\beta:\Lambda^2V\to W.
\]

Conversely, every alternating tensor defines a class-two exponent-$p$ group

\[
G_\beta=V\oplus W
\]

with multiplication

\[
(v,z)(v',z')
\left(v+v',z+z'+\frac12\beta(v,v')\right),
\]

and commutator

\[
[(v,z),(v',z')]=(0,\beta(v,v')).
\]

Two such special groups are isomorphic if and only if their commutator tensors are pseudo-isometric. Thus, reconstructing

\[
\beta:\Lambda^2V\to W
\]

up to the action of

\[
\operatorname{GL}(V)\times \operatorname{GL}(W)
\]

is a natural form of the class-two exponent-$p$ group isomorphism problem.

For $\dim W=2$, the tensor is an alternating pencil. The complete Kronecker canonical data of such pencils can be recovered from explicitly designed commutator-word distributions. The present paper addresses the first genuinely higher-dimensional reconstruction mechanism.

Assume

\[
\dim V=2r+1
\]

is odd. For every functional

\[
\lambda\in W^*
\]

we have an alternating form

\[
A_\lambda=\lambda\circ\beta.
\]

Because $\dim V$ is odd, each $A_\lambda$ is singular. The generic corank is at least (1). The first nontrivial degeneracy condition is

\[
\dim\ker A_\lambda\ge3.
\]

This condition is cut out by the submaximal Pfaffians of the universal alternating matrix

\[
M_\beta(x)=x_1A_1+\cdots+x_mA_m.
\]

The central idea is:

\[
\boxed{
\text{MVMT rank data recovers the Pfaffian corank-}(\ge3)\text{ scheme.}
}
\]

On a generic codimension-three Pfaffian locus, the Buchsbaum--Eisenbud theorem says that the ideal is governed by a skew-symmetric matrix in the minimal free resolution. Thus, under suitable rigidity, recovering the Pfaffian ideal recovers the tensor.

The result is deliberately formulated with explicit hypotheses. The reduced degeneracy scheme determines the Pfaffian ideal only under radicality and saturation. The Pfaffian ideal determines the matrix through the Buchsbaum--Eisenbud resolution. Left-right recovery of the matrix becomes congruence only under a self-adjoint rigidity hypothesis. These hypotheses hold on a nonempty Zariski open set in the range

\[
4\le \dim W\le10.
\]

\subsection{Main theorem}

Let

\[
n=2r+1\ge5,\qquad 4\le m\le10.
\]

Let

\[
\mathcal T_{n,m}=\operatorname{Hom}(\Lambda^2V,W)
\]

with $\dim V=n$, $\dim W=m$.

\begin{theorem}[Theorem A -- Generic odd-Pfaffian MVMT reconstruction]

There exists a nonempty Zariski open subset

\[
\mathcal U_{\mathrm{rec}}\subset \mathcal T_{n,m}
\]

such that, for every

\[
\beta\in\mathcal U_{\mathrm{rec}},
\]

the pointwise Fourier-tagged scalar-extension commutator-word moment data determine $\beta$ up to pseudo-isometry.

Equivalently, for every surjective radical-free

\[
\beta\in\mathcal U_{\mathrm{rec}},
\]

the associated special class-two exponent-$p$ group $G_\beta$ is determined up to isomorphism by these word-moment data.

\end{theorem}

\subsection{What ``pointwise Fourier-tagged'' means}

The theorem uses more than aggregate rank-zeta data. It uses the Fourier coefficient attached to each

\[
\lambda\in W_K^*
\]

over every finite extension

\[
K/k.
\]

That coefficient gives

\[
|K|^{-\operatorname{rank}(\lambda\circ\beta_K)}.
\]

Thus the data determine the rank pointwise on $W_K^*$, not merely the number of points of each rank.

\subsection{What is not claimed}

We do not claim:

\begin{itemize}
\item classification of all alternating tensors;
\end{itemize}

\begin{itemize}
\item recovery of nonreduced Pfaffian schemes from reduced rank data alone;
\end{itemize}

\begin{itemize}
\item an optimized group-isomorphism algorithm;
\end{itemize}

\begin{itemize}
\item arbitrary target dimension;
\end{itemize}

\begin{itemize}
\item completeness from rank-zeta point counts alone.
\end{itemize}

The theorem is a generic reconstruction theorem on an explicitly described open locus.

\section{Alternating tensors and the universal odd Pfaffian matrix}

Let

\[
k=\mathbb F_p,\qquad p\neq2.
\]

Let

\[
V
\]

be a vector space of dimension

\[
n=2r+1\ge5.
\]

Let

\[
W
\]

be a vector space of dimension $m$, and let

\[
\beta:\Lambda^2 V\to W
\]

be alternating.

Choose a basis

\[
w_1,\ldots,w_m
\]

of $W$, and let

\[
x_1,\ldots,x_m
\]

be the dual homogeneous coordinates on $W^*$. Write

\[
\beta=(A_1,\ldots,A_m),
\qquad
A_i\in\Lambda^2V^*.
\]

Set

\[
S=k[x_1,\ldots,x_m].
\]

Define the universal alternating matrix

\[
M_\beta(x)=x_1A_1+\cdots+x_mA_m.
\]

This is an $n\times n$ skew-symmetric matrix with linear entries in $S$.

Since $n=2r+1$ is odd,

\[
\det M_\beta(x)=0
\]

identically. The submaximal Pfaffians, i.e. the Pfaffians of the $2r\times2r$ principal submatrices obtained by deleting one row and column, are homogeneous forms of degree $r$.

Let

\[
p_1,\ldots,p_n
\]

be the signed submaximal Pfaffians of $M_\beta$, and define

\[
I_\beta=(p_1,\ldots,p_n)\subset S.
\]

Define the projective Pfaffian degeneracy scheme

\[
Z_\beta=\operatorname{Proj}(S/I_\beta)
\subseteq \mathbb P(W^*).
\]

Set-theoretically,

\[
Z_\beta
=
\{[\lambda]\in\mathbb P(W^*):\dim\ker(\lambda\circ\beta)\ge3\}.
\]

\section{Pointwise Fourier-tagged MVMT recovers the reduced Pfaffian locus}

Let

\[
K/k
\]

be a finite extension. Extend

\[
\beta
\]

to

\[
\beta_K:\Lambda^2_KV_K\to W_K.
\]

Let

\[
\mu_{\beta,K}
\]

be the distribution of the commutator output

\[
\beta_K(u,v)\in W_K
\]

where

\[
u,v\in V_K
\]

are uniform.

For

\[
\lambda\in W_K^*
\]

let

\[
A_\lambda=\lambda\circ\beta_K.
\]

Let

\[
\chi:K\to\mathbb C^\times
\]

be a nontrivial additive character. The Fourier coefficient of $\mu_{\beta,K}$ at $\lambda$ is

\[
\widehat\mu_{\beta,K}(\lambda)
=
\mathbb E_{u,v}
\chi(A_\lambda(u,v)).
\]

\begin{lemma}[Pointwise rank recovery]

For every finite extension $K/k$ and every

\[
\lambda\in W_K^*,
\]

one has

\[
\boxed{
\widehat\mu_{\beta,K}(\lambda)
=
|K|^{-\operatorname{rank}_K(A_\lambda)}.
}
\]

Therefore the Fourier-tagged scalar-extension MVMT data determine

\[
\operatorname{rank}_K(A_\lambda)
\]

for every $\lambda\in W_K^*$.

\end{lemma}

\begin{proof}

This is the standard bilinear character average. For fixed $u$, the average over $v$ is $1$ precisely when

\[
A_\lambda(u,-)=0,
\]

and $0$ otherwise. Hence the total average is the proportion of $u$ in the radical of $A_\lambda$, namely

\[
|K|^{-\operatorname{rank}(A_\lambda)}.
\]

\end{proof}

\begin{lemma}[MVMT recovers $I_\beta$ under radicality and saturation]

Assume that

\[
I_\beta
\]

is saturated and geometrically radical. Then the pointwise Fourier-tagged scalar-extension MVMT data determine

\[
I_\beta.
\]

\end{lemma}

\begin{proof}

By Lemma 3.1, MVMT determines the rank of $A_\lambda$ for every

\[
\lambda\in W_K^*
\]

over every finite extension $K/k$. Therefore it determines, for every $K/k$,

\[
Z_\beta(K)
=
\{[\lambda]\in\mathbb P(W_K^*):\operatorname{rank}(A_\lambda)\le2r-2\}.
\]

Equivalently, it determines the $K$-rational points of the corank-$\ge3$ locus for every finite extension $K$.

Every closed point of a finite-type $k$-scheme is rational over some finite extension of $k$. Thus the MVMT data determine the set of all closed points of

\[
(Z_\beta)_{\mathrm{red}}.
\]

Since $I_\beta$ is geometrically radical and saturated, it is exactly the homogeneous ideal of this reduced projective subscheme. Hence $I_\beta$ is determined.

\end{proof}

\begin{remark}[Remark 3.3]

The radicality assumption is essential. Pointwise rank data over finite fields detects the reduced degeneracy locus. Without an Artin or scheme-theoretic enhancement, it need not determine nilpotent structure in $I_\beta$.

\end{remark}

\section{Buchsbaum--Eisenbud recovery from the Pfaffian ideal}

We now recall the codimension-three Pfaffian structure.

Assume

\[
\operatorname{ht}I_\beta=3.
\]

For an odd skew-symmetric matrix

\[
M_\beta
\]

of size $n=2r+1$, the Buchsbaum--Eisenbud Pfaffian complex has the form

\[
0
\longrightarrow
S(-2r-1)
\xrightarrow{p^{\mathsf T}}
S(-r-1)^n
\xrightarrow{M_\beta}
S(-r)^n
\xrightarrow{p}
I_\beta
\longrightarrow0,
\]

where

\[
p=(p_1,\ldots,p_n)
\]

is the vector of signed submaximal Pfaffians.

The Buchsbaum--Eisenbud theorem says that codimension-three Gorenstein ideals are governed by such alternating matrices. In our setting, we use the following precise consequence.

\begin{lemma}[The Pfaffian ideal determines the middle matrix up to left-right equivalence]

Assume:

\begin{itemize}
\item $I_\beta$ has height $3$;
\end{itemize}

\begin{itemize}
\item the Buchsbaum--Eisenbud Pfaffian complex above is the minimal graded free resolution of $I_\beta$.
\end{itemize}

Let

\[
M_{\beta'}
\]

be another odd skew-symmetric linear matrix such that

\[
I_{\beta'}=I_\beta
\]

and such that its Buchsbaum--Eisenbud Pfaffian complex is also the minimal graded free resolution of $I_\beta$.

Then there exist

\[
P,Q\in \operatorname{GL}_n(k)
\]

such that

\[
M_{\beta'}=P\,M_\beta\,Q.
\]

\end{lemma}

\begin{proof}

Both matrices give minimal graded free resolutions of the same homogeneous ideal $I_\beta$:

\[
0
\to
S(-2r-1)
\to
S(-r-1)^n
\xrightarrow{M_\beta}
S(-r)^n
\to
I_\beta
\to0,
\]

and

\[
0
\to
S(-2r-1)
\to
S(-r-1)^n
\xrightarrow{M_{\beta'}}
S(-r)^n
\to
I_\beta
\to0.
\]

Minimal graded free resolutions over $S$ are unique up to graded chain isomorphism. Since

\[
S(-r-1)^n
\]

and

\[
S(-r)^n
\]

are pure sums of equal shifts, every graded automorphism of these middle terms is given by a constant invertible $n\times n$ matrix over $k$. Therefore the middle maps differ by constant invertible transformations:

\[
M_{\beta'}=P\,M_\beta\,Q
\]

for some

\[
P,Q\in \operatorname{GL}_n(k).
\]

\end{proof}

\begin{remark}[Remark 4.2]

This lemma gives only left-right equivalence of the matrices. Pseudo-isometry of alternating tensors requires congruence. The next section gives the additional rigidity condition needed to pass from left-right equivalence to congruence.

\end{remark}

\section{Self-adjoint rigidity}

For

\[
\beta=(A_1,\ldots,A_m),
\]

define

\[
\operatorname{Adj}^+(\beta)
=
\{R\in\operatorname{End}(V):R^{\mathsf T}A_i=A_iR\text{ for every }i\}.
\]

Equivalently, for

\[
M_\beta(x)=\sum_i x_iA_i,
\]

we have

\[
R\in\operatorname{Adj}^+(\beta)
\]

if and only if

\[
R^{\mathsf T}M_\beta=M_\beta R.
\]

\begin{lemma}[Rigidity upgrades left-right equivalence to congruence]

Let

\[
M=M_\beta
\]

and

\[
M'=M_{\beta'}
\]

be alternating linear matrices of the same size. Suppose

\[
M'=P M Q
\]

for some

\[
P,Q\in \operatorname{GL}(V).
\]

Assume

\[
\operatorname{Adj}^+(\beta)=kI.
\]

Then

\[
M'=c P M P^{\mathsf T}
\]

for some

\[
c\in k^\times.
\]

Consequently, $M$ and $M'$ define pseudo-isometric alternating tensors.

\end{lemma}

\begin{proof}

Since both $M$ and $M'$ are alternating,

\[
M^{\mathsf T}=-M,
\qquad
(M')^{\mathsf T}=-M'.
\]

From

\[
M'=PMQ
\]

we get

\[
(M')^{\mathsf T}
=
Q^{\mathsf T}M^{\mathsf T}P^{\mathsf T}
=
-Q^{\mathsf T}MP^{\mathsf T}.
\]

But

\[
(M')^{\mathsf T}=-M'=-PMQ.
\]

Hence

\[
Q^{\mathsf T}MP^{\mathsf T}=PMQ.
\]

Multiplying on the left by $P^{-1}$ and on the right by $P^{-\mathsf T}$, we obtain

\[
P^{-1}Q^{\mathsf T}M
=
M QP^{-\mathsf T}.
\]

Set

\[
R=QP^{-\mathsf T}.
\]

Then

\[
R^{\mathsf T}=P^{-1}Q^{\mathsf T}.
\]

Therefore

\[
R^{\mathsf T}M=MR.
\]

Since

\[
M=\sum_i x_iA_i,
\]

this identity is equivalent to

\[
R^{\mathsf T}A_i=A_iR
\qquad
\text{for every }i.
\]

Thus

\[
R\in\operatorname{Adj}^+(\beta).
\]

By hypothesis,

\[
R=cI
\]

for some

\[
c\in k^\times.
\]

Therefore

\[
Q=cP^{\mathsf T}.
\]

Substituting into

\[
M'=PMQ
\]

gives

\[
M'=c P M P^{\mathsf T}.
\]

The scalar $c$ is absorbed by a scalar change of target coordinates in $W$. Hence the corresponding tensors are pseudo-isometric.

\end{proof}

\section{Genericity of the Pfaffian hypotheses}

We now show that the hypotheses used above hold on a nonempty Zariski open subset in the range

\[
4\le m\le10.
\]

\begin{lemma}[Generic Pfaffian degeneracy hypotheses]

Assume

\[
n=2r+1\ge5,
\qquad
4\le m\le10.
\]

There exists a nonempty Zariski open subset

\[
\mathcal U_{\mathrm{Pf}}
=
\subset
\operatorname{Hom}(\Lambda^2V,W)
\]

such that, for every

\[
\beta\in\mathcal U_{\mathrm{Pf}},
\]

the Pfaffian ideal $I_\beta$ has height (3), is saturated and geometrically radical, and the Buchsbaum--Eisenbud Pfaffian complex is its minimal free resolution.

\end{lemma}

\begin{proof}

Let

\[
\mathbb P(\Lambda^2V^*)
\]

be the projective space of alternating forms on $V$.

Let

\[
\Sigma_3
=
\{[\omega]:\dim\ker\omega\ge3\}
\]

be the corank-$\ge3$ Pfaffian variety. Since

\[
\dim V=2r+1
\]

is odd, $\Sigma_3$ is the first nontrivial degeneracy locus. It has codimension $3$. Its singular locus is

\[
\Sigma_5
=
\{[\omega]:\dim\ker\omega\ge5\},
\]

which has codimension $10$.

A tensor

\[
\beta:\Lambda^2V\to W
\]

is the same as a projective linear subspace

\[
\mathbb P(W^*)\subset\mathbb P(\Lambda^2V^*).
\]

For a general such $\mathbb P(W^*)$, the intersection

\[
Z_\beta=\mathbb P(W^*)\cap\Sigma_3
\]

has expected codimension $3$ in $\mathbb P(W^*)$. Since

\[
m-1<10
\]

for

\[
m\le10,
\]

a general $\mathbb P(W^*)$ avoids $\Sigma_5$. Therefore the intersection occurs in the smooth locus of $\Sigma_3$.

By generic transversality, a general $\mathbb P(W^*)$ intersects the smooth locus of $\Sigma_3$ transversely. Hence $Z_\beta$ is smooth over the algebraic closure. In particular, its homogeneous ideal is saturated and geometrically radical.

The ideal of $\Sigma_3$ is generated by the submaximal Pfaffians of the universal odd alternating matrix. The ideal $I_\beta$ is the pullback of this ideal to $\mathbb P(W^*)$. Since the intersection has expected codimension $3$, the Buchsbaum--Eisenbud acyclicity criterion gives exactness of the associated Pfaffian complex. Minimality follows because all entries of $M_\beta$ are linear and have no constant terms, and the maps are homogeneous of positive degree.

The conditions above are open conditions on $\beta$, and the general-position argument shows that the open set is nonempty.

\end{proof}

\section{Genericity of self-adjoint rigidity}

We now prove that

\[
\operatorname{Adj}^+(\beta)=kI
\]

is also a nonempty open condition.

\begin{lemma}[Self-adjoint rigidity is open]

The condition

\[
\dim_k\operatorname{Adj}^+(\beta)=1
\]

is Zariski open in

\[
\operatorname{Hom}(\Lambda^2V,W).
\]

\end{lemma}

\begin{proof}

For fixed $\beta=(A_1,\ldots,A_m)$, the equations

\[
R^{\mathsf T}A_i=A_iR
\]

are linear equations in the entries of $R$, with coefficients depending linearly on the entries of the forms $A_i$. Thus

\[
\operatorname{Adj}^+(\beta)
\]

is the kernel of a matrix whose entries are polynomial functions of $\beta$. Kernel dimension is upper semicontinuous. Since scalar matrices always belong to $\operatorname{Adj}^+(\beta)$, the condition that the dimension be exactly $1$ is open.

\end{proof}

It remains to show nonemptiness. We give an explicit pair of alternating forms whose common self-adjoint algebra is scalar.

Let

\[
V=X\oplus Y,
\]

where

\[
X=\langle x_0,\ldots,x_r\rangle,
\qquad
Y=\langle y_0,\ldots,y_{r-1}\rangle.
\]

Define

\[
A_0=\sum_{i=0}^{r-1}x_i^*\wedge y_i^*,
\]

\[
A_1=\sum_{i=0}^{r-1}x_{i+1}^*\wedge y_i^*.
\]

These are the two alternating forms of the odd Kronecker block.

\begin{lemma}[The odd Kronecker pair has scalar self-adjoint algebra]

For the pair $(A_0,A_1)$ above,

\[
\operatorname{Adj}^+(A_0,A_1)=kI_V.
\]

\end{lemma}

\begin{proof}

Write an endomorphism

\[
T\in\operatorname{End}(V)
\]

in block form relative to

\[
V=X\oplus Y:
\]

\[
T=
\begin{pmatrix}
P&Q\
C&S
\end{pmatrix},
\]

where

\[
P:X\to X,\qquad Q:Y\to X,\qquad C:X\to Y,\qquad S:Y\to Y.
\]

Let

\[
L_0=[I_r\ 0],
\qquad
L_1=[0\ I_r],
\]

viewed as $r\times(r+1)$ matrices. The forms $A_0,A_1$ have block matrices

\[
\begin{pmatrix}
0&L_j^{\mathsf T}\
-L_j&0
\end{pmatrix},
\qquad
j=0,1.
\]

The condition

\[
T^{\mathsf T}A_j=A_jT
\]

is equivalent to the block equations

\[
S^{\mathsf T}L_j=L_jP,
\tag{1}
\]

\[
L_jQ+(L_jQ)^{\mathsf T}=0,
\tag{2}
\]

\[
L_j^{\mathsf T}C+C^{\mathsf T}L_j=0.
\tag{3}
\]

First consider (1). From $j=0$, since

\[
L_0=[I_r\ 0],
\]

the first $r$ rows of $P$ are

\[
[S^{\mathsf T}\ 0].
\]

From $j=1$, since

\[
L_1=[0\ I_r],
\]

the last $r$ rows of $P$ are

\[
[0\ S^{\mathsf T}].
\]

Comparing the overlapping rows of $P$ gives

\[
S_{a,b}=S_{a-1,b-1}
\]

for

\[
1\le a,b\le r-1,
\]

and the boundary equations give

\[
S_{0,b}=0\quad(b>0),
\qquad
S_{a,0}=0\quad(a>0).
\]

Thus all off-diagonal entries of $S$ vanish and all diagonal entries are equal. Hence

\[
S=cI_Y.
\]

The same equations then force

\[
P=cI_X.
\]

Now consider $Q=(q_{ij})$, with

\[
0\le i\le r,
\qquad
0\le j\le r-1.
\]

Equation (2) for $j=0$ gives

\[
q_{ij}+q_{ji}=0
\qquad
(0\le i,j\le r-1).
\tag{4}
\]

Equation (2) for $j=1$ gives

\[
q_{i+1,j}+q_{j+1,i}=0
\qquad
(0\le i,j\le r-1).
\tag{5}
\]

From (4), $q_{ii}=0$. From (5), $q_{i+1,i}=0$. If $a<b$, then applying (4) and then (5) gives

\[
q_{a,b}
-q_{b,a}
q_{a+1,b-1}.
\]

This reduces (b-a) by (2), eventually reaching a diagonal or subdiagonal entry, both of which are zero. Hence $q_{a,b}=0$ for $a<b$. The case $a>b$ is similar, using (5) and (4) to reduce to a diagonal or subdiagonal zero. Therefore

\[
Q=0.
\]

The equations for $C$ are the transpose-dual of the equations for $Q$, and the same argument gives

\[
C=0.
\]

Thus

\[
T=
\begin{pmatrix}
cI_X&0\
0&cI_Y
\end{pmatrix}
cI_V.
\]

Therefore

\[
\operatorname{Adj}^+(A_0,A_1)=kI_V.
\]

\end{proof}

\begin{corollary}[Nonempty open self-adjoint locus]

For every

\[
m\ge2,
\]

the locus

\[
\operatorname{Adj}^+(\beta)=kI
\]

is a nonempty Zariski open subset of

\[
\operatorname{Hom}(\Lambda^2V,W).
\]

\end{corollary}

\begin{proof}

By Lemma 7.1 it is open. By Lemma 7.2, the pair $(A_0,A_1)$ has scalar self-adjoint algebra. Adding more alternating forms can only decrease the common self-adjoint algebra. Hence the condition is nonempty for all $m\ge2$.

\end{proof}

\section{The generic reconstruction locus}

Let

\[
\mathcal U_{\mathrm{Pf}}
\]

be the open set from Lemma 6.1, and let

\[
\mathcal U_{\mathrm{Adj}}
=
\{\beta:\operatorname{Adj}^+(\beta)=kI\}.
\]

Define

\[
\mathcal U_{\mathrm{rec}}
=
\mathcal U_{\mathrm{Pf}}\cap \mathcal U_{\mathrm{Adj}}.
\]

Since

\[
\operatorname{Hom}(\Lambda^2V,W)
\]

is irreducible and both open sets are nonempty, their intersection is nonempty.

\begin{theorem}[Generic odd-Pfaffian MVMT reconstruction]

Let

\[
k=\mathbb F_p,\qquad p\neq2,
\]

let

\[
\dim V=2r+1\ge5,
\]

and let

\[
4\le \dim W=m\le10.
\]

For every

\[
\beta\in\mathcal U_{\mathrm{rec}},
\]

the pointwise Fourier-tagged scalar-extension MVMT data determine $\beta$ up to pseudo-isometry.

\end{theorem}

\begin{proof}

By Lemma 3.2, the MVMT data determine the Pfaffian ideal

\[
I_\beta.
\]

By Lemma 4.1, this ideal determines the universal alternating matrix

\[
M_\beta
\]

up to left-right equivalence. By Lemma 5.1 and the condition

\[
\operatorname{Adj}^+(\beta)=kI,
\]

this left-right equivalence upgrades to congruence up to target scalar, hence to pseudo-isometry of the alternating tensor.

\end{proof}

\section{Group-theoretic corollary}

Suppose additionally that

\[
\beta:\Lambda^2V\to W
\]

is surjective and radical-free:

\[
\operatorname{im}\beta=W,
\qquad
\operatorname{rad}\beta=0.
\]

Then $G_\beta$ is a special class-two exponent-$p$ group with

\[
G_\beta'=Z(G_\beta)=W.
\]

\begin{corollary}[Generic group reconstruction]

For every surjective radical-free

\[
\beta\in\mathcal U_{\mathrm{rec}},
\]

the group

\[
G_\beta
\]

is determined up to isomorphism by the pointwise Fourier-tagged scalar-extension MVMT data.

\end{corollary}

\begin{proof}

For special class-two exponent-$p$ groups, group isomorphism is equivalent to pseudo-isometry of the commutator tensor. Theorem 8.1 recovers the pseudo-isometry class of $\beta$.

\end{proof}

\section{Density over large finite fields}

The theorem above is geometric. It gives a nonempty Zariski open subset. Over finite fields, one may translate this into a density statement after spreading out.

\begin{corollary}[Density-one form]

Fix

\[
n=2r+1\ge5,
\qquad
4\le m\le10.
\]

After choosing an integral model of the open set

\[
\mathcal U_{\mathrm{rec}},
\]

for all sufficiently large finite fields $\mathbb F_q$ of odd characteristic,

\[
\frac{
|\mathcal U_{\mathrm{rec}}(\mathbb F_q)|
}{
|\operatorname{Hom}(\Lambda^2\mathbb F_q^n,\mathbb F_q^m)|
}
1-O_{n,m}(q^{-1}).
\]

\end{corollary}

\begin{proof}

The complement of a nonempty Zariski open subset in affine space is contained in a proper algebraic subvariety of bounded degree after spreading out over a finitely generated base. The number of $\mathbb F_q$-points of this complement is

\[
O(q^{D-1}),
\]

where

\[
D=\dim \operatorname{Hom}(\Lambda^2 k^n,k^m).
\]

Dividing by $q^D$ gives

\[
O(q^{-1}).
\]

\end{proof}

\section{Scope and limitations}

The theorem proves generic reconstruction beyond alternating pencils, but it is not a universal theorem.

It does not claim:

\begin{itemize}
\item recovery of nonreduced Pfaffian schemes from rank data;
\end{itemize}

\begin{itemize}
\item classification of all alternating tensors;
\end{itemize}

\begin{itemize}
\item arbitrary target dimension;
\end{itemize}

\begin{itemize}
\item polynomial-time isomorphism testing;
\end{itemize}

\begin{itemize}
\item classification from aggregate rank-zeta point counts.
\end{itemize}

The theorem does claim:

\[
\boxed{
\text{pointwise Fourier-tagged MVMT recovers the reduced Pfaffian degeneracy scheme;}
}
\]

\[
\boxed{
\text{on a generic rigid codimension-three Pfaffian open locus, this scheme determines the tensor;}
}
\]

\[
\boxed{
\text{therefore MVMT is generically complete beyond the pencil case.}
}
\]

\begin{thebibliography}{99}
\bibitem{ref1} D. A. Buchsbaum and D. Eisenbud, Algebra structures for finite free resolutions, and some structure theorems for ideals of codimension 3, American Journal of Mathematics 99 (1977), 447--485.
\bibitem{ref2} D. Eisenbud, S. Popescu, and C. Walter, Lagrangian subbundles and codimension 3 subcanonical subschemes, arXiv:math/9906170.
\bibitem{ref3} A. Beauville, Determinantal hypersurfaces, Michigan Math. J. 48 (2000), 39--64.
\bibitem{ref4} D. Kerner and V. Vinnikov, Determinantal representations of singular hypersurfaces in $\mathbb P^n$, arXiv:0906.3012.
\bibitem{ref5} M. Levy, Enumerating fibres of commutator words over $p$-groups, arXiv:1602.04093.
\bibitem{ref6} J. B. Wilson, Decomposing $p$-groups via Jordan algebras, Journal of Algebra 322 (2009), 2642--2679.
\bibitem{ref7} P. A. Brooksbank and J. B. Wilson, Computing isometry groups of Hermitian maps, Transactions of the American Mathematical Society 364 (2012), 1975--1996.
\bibitem{ref8} X. Sun, Faster Isomorphism for $p$-Groups of Class 2 and Exponent $p$, arXiv:2303.15412.
\end{thebibliography}
\end{document}
